% Proposed insertion after the local contraction and strict stalling results.
% Uses only the existing amsmath/amsthm/enumitem setup and shared macros.
\begin{example}[A critical tangent map with sublinear tightening]
\label{ex:critical-scalar}
Let $X=\R$, $f(x)=x^2$, and $B_0=[-1/2,1/2]$, so that
$x^*=0$ and $f^*=0$. For every nonempty subbox $B=[\ell,u]\subseteq B_0$,
including degenerate intervals, define
\begin{equation}\label{eq:critical-relaxation}
  E(w)=\frac{w^2}{4}-\frac{w^3}{8},\qquad
  \phi_B:B\longrightarrow\R\subset\R\cup\{+\infty\},\qquad
  \phi_B(x)=x^2-E(w(B)).
\end{equation}
This is a valid monotone family of continuous convex projected objectives.
Its cutoff sets are compact, and we use $T_U(\emptyset)=\emptyset$.
At $U=0$, no positive-width box is fixed; the only nonempty fixed box is
$\{0\}$. Nevertheless, the following properties hold.
\begin{enumerate}[label=(\alph*)]
\item For $d=(d^-,d^+)\ge0$ and $D(d)=[-d^-,d^+]$, the tangent model is
  \begin{equation}\label{eq:critical-tangent}
    Q(d,\xi)=\xi^2-\frac{(d^-+d^+)^2}{4}.
  \end{equation}
  The remainder is $O(t^3)$ uniformly on compact shape sets, $r^*=1$,
  and the two face values at $d=(1,1)$ are zero. Thus neither strict face
  condition holds at this shape.
\item For a symmetric initial box $[-h_0,h_0]$ with $0<h_0\le1/2$, the
  zero-cutoff Jacobi radii satisfy
  \begin{equation}\label{eq:critical-recurrence}
    h_{k+1}=h_k\sqrt{1-h_k},\qquad h_k\sim\frac{2}{k}.
  \end{equation}
\item For the same initial box and cutoff $U=\epsilon$ with
  $0\le\epsilon\le h_0^3$, the limiting radius is
  $h_\infty=\epsilon^{1/3}$ and the limiting width is
  $2\epsilon^{1/3}$. If $0<\epsilon<h_0^3$, the radii decrease strictly
  to this limit; if $\epsilon=h_0^3$, the initial box is fixed.
\end{enumerate}
The family is an abstract projected relaxation. It demonstrates that the
local contraction and strict stalling criteria are sufficient conditions
and do not exhaust the possible behavior of a valid monotone family.
\end{example}

\begin{proof}
Every subbox has $0\le w(B)\le1$. On this interval $E(w)\ge0$ and
\[
  E'(w)=\frac{w}{2}-\frac{3w^2}{8}\ge0.
\]
Consequently $\phi_B(x)\le f(x)$ for every $x\in X\cap B$. If
$B'\subseteq B$, then $w(B')\le w(B)$, so for every $x\in B'$,
\[
  \phi_{B'}(x)=x^2-E(w(B'))\ge x^2-E(w(B))=\phi_B(x).
\]
Continuity and convexity hold on each box, and compactness of the sublevel
sets follows from continuity and boundedness of $B$. For $U\ge0$ their
exact box hull is the interval, possibly empty,
\begin{equation}\label{eq:critical-operator}
  T_U(B)=B\cap
  [-\sqrt{U+E(w(B))},\sqrt{U+E(w(B))}].
\end{equation}
If $w=w(B)>0$, this interval at $U=0$ is contained in an interval of width
\[
  2\sqrt{E(w)}=w\sqrt{1-w/2}<w.
\]
It therefore cannot equal $B$, whether or not $B$ contains zero or is
symmetric. On a singleton $\{x\}$, the projected value is $x^2$; hence
$\{0\}$ is the only nonempty fixed box at cutoff zero.

For the tangent statement, write $s=d^-+d^+$. Whenever $tD(d)\subseteq B_0$,
the exact identity is
\[
  \phi_{tD(d)}(t\xi)
  =t^2\left(\xi^2-\frac{s^2}{4}\right)+\frac{t^3s^3}{8}.
\]
The cubic term divided by $t^2$ tends to zero uniformly for $d$ in any
compact subset of $\R^2_{\ge0}$. For $c\ge0$ put
$q_c(d)=\sqrt{s^2/4+c}$. Taking the hull of the tangent sublevel set gives
\[
  \Phi_c(d)=
  \bigl(\min\{d^-,q_c(d)\},\min\{d^+,q_c(d)\}\bigr).
\]
For every $d\gg0$, its smaller component is at most $s/2\le q_c(d)$
and is therefore unchanged by $\Phi_c$. No inequality
$\Phi_\delta(d)\le\lambda d$ with $\delta>0$ can hold with $\lambda<1$.
The choice $\lambda=1$ always works, so the defining infimum gives
$r^*=1$. At $d=(1,1)$ one has $Q(d,\xi)=\xi^2-1$, whose values at both
faces $\xi=\pm1$ are zero. These are non-strict tangent face witnesses,
but the actual relaxed face values on $[-t,t]$ are
$\phi_{[-t,t]}(\pm t)=t^3>0$. Thus replacing strict negative face slack
by non-strict witnesses would give a false stalling conclusion here.

By~\eqref{eq:critical-operator}, a symmetric box remains symmetric. At
cutoff zero its new radius is $\sqrt{E(2h)}=h\sqrt{1-h}$, proving the
recurrence. The positive radii decrease to a limit $H\ge0$. Continuity
gives $H=H\sqrt{1-H}$, hence $H=0$. Moreover,
\[
  \frac1{h_{k+1}}-\frac1{h_k}
  =\frac1{\sqrt{1-h_k}(1+\sqrt{1-h_k})}\longrightarrow\frac12.
\]
Summing these increments and dividing by $k$ gives
$(1/h_k)/k\to1/2$, which proves $h_k\sim2/k$.

At cutoff $\epsilon\ge0$, the exact radius update is
\[
  h_{k+1}=\min\{h_k,g_\epsilon(h_k)\},\qquad
  g_\epsilon(h)=\sqrt{\epsilon+h^2-h^3}.
\]
The case $\epsilon=0$ was just proved. For $\epsilon>0$, put
$a=\epsilon^{1/3}\le h_0\le1/2$. Then $g_\epsilon(a)=a$, and
$g_\epsilon$ is strictly increasing on $(0,1/2]$, because the derivative
of its squared value is $2h-3h^2>0$. For $a<h\le1/2$ it follows that
$a<g_\epsilon(h)<h$, the second inequality using $\epsilon<h^3$.
If $a<h_0$, the radii therefore decrease strictly to a limit $H\ge a$.
Continuity gives $H=g_\epsilon(H)$, so $H^3=\epsilon$ and $H=a$.
If $a=h_0$, the radius stays equal to $a$. Both cases give the asserted
limit, including the terminal equality $\epsilon=h_0^3$.
\end{proof}
